% Generated draft: requires mathematical and source review.
\documentclass[11pt]{article}
\usepackage[margin=25mm]{geometry}
\usepackage{fontspec}
\setmainfont{Latin Modern Roman}
\usepackage{amsmath,amssymb,mathtools,unicode-math}
\setmathfont{Latin Modern Math}
\usepackage{enumitem,needspace,xurl,hyperref}
\setlength{\parindent}{0pt}
\setlength{\parskip}{5pt}
\setlength{\emergencystretch}{3em}
\newcommand{\N}{\mathbb{N}}
\newcommand{\Z}{\mathbb{Z}}
\newcommand{\Q}{\mathbb{Q}}
\newcommand{\R}{\mathbb{R}}
\newcommand{\C}{\mathbb{C}}
\newcommand{\F}{\mathbb{F}}
\newcommand{\A}{\mathbb{A}}
\newcommand{\PP}{\mathbb{P}}
\newcommand{\E}{\mathbb{E}}
\newcommand{\eps}{\varepsilon}
\newcommand{\dd}{\,\mathrm{d}}
\begin{document}
\section*{UnsolvedMath ID: 30005320}
\textbf{Source problem code:} catalogue entry
\subsection{Definitions and mathematical statement}
Under the standard classical-pattern interpretation, let \(S_n\) be the permutations of \(\{1,\ldots,n\}\), with \(S_0\) consisting of the empty permutation. For \(\tau\in S_m\) and \(\pi\in S_n\), define
\[
\operatorname{occ}_{\tau}(\pi)=\#\{(i_1,\ldots,i_m):1\le i_1<\cdots<i_m\le n,\ \pi_{i_1}\cdots\pi_{i_m}\text{ has relative order }\tau\}.
\]
For integers \(r,n\ge0\), put
\[
\psi_{\tau,r}(n)=\#\{\pi\in S_n:\operatorname{occ}_{\tau}(\pi)=r\},
\qquad
\operatorname{Av}_n(\tau)=\{\pi\in S_n:\operatorname{occ}_{\tau}(\pi)=0\}.
\]
Thus \(\psi_{\tau,0}(n)=|\operatorname{Av}_n(\tau)|\). Two patterns are Wilf-equivalent when their avoidance counts agree for every \(n\). This definition alone does not identify their counts with a positive prescribed number of occurrences. The supplied record does not explicitly specify the occurrence convention; the definitions above make the standard interpretation explicit.

The asymptotic question concerns two unspecified Wilf classes of patterns of length four. For a pattern \(\tau\) representing either intended class and each fixed integer \(r\ge0\), the supplied conjectural formula is
\[
\psi_{\tau,r}(n)\sim C_{\tau,r}\,n^r\psi_{\tau,0}(n)
\qquad(n\to\infty).
\]
Here the usual positive-constant interpretation would require a constant \(C_{\tau,r}>0\), independent of \(n\), such that
\[
\lim_{n\to\infty}\frac{\psi_{\tau,r}(n)}{n^r\psi_{\tau,0}(n)}=C_{\tau,r}.
\]
The source writes \(C_r\); the additional pattern subscript records dependence that cannot be excluded from the supplied evidence. The two classes, their representatives, the intended range of \(r\), and any invariance of positive-occurrence counts within a class are not supplied. Consequently, this is a conditional reconstruction of the intended question, not a conjecture asserted here for every pattern of length four.

For the moment-sequence component, define the formal series
\[
F_{\tau,r}(x)=\sum_{n\ge0}\psi_{\tau,r}(n)x^n,
\qquad
g_{\tau,n}(q)=\sum_{k=0}^{\binom{n}{m}}q^k\psi_{\tau,k}(n),
\qquad
G_{\tau}(x;q)=\sum_{n\ge0}g_{\tau,n}(q)x^n.
\]
Here \(\binom{n}{m}=0\) for \(n<m\). The expression for \(g_{\tau,n}(q)\) is the natural reconstruction of the damaged displayed sum in the record; the variable \(x\) has been introduced to distinguish the ordinary generating-function variable from the occurrence weight \(q\).

A real sequence \((a_n)_{n\ge0}\) is a Stieltjes moment sequence if there is a nonnegative Borel measure \(\mu\) on \([0,\infty)\), with every moment finite, such that
\[
a_n=\int_{[0,\infty)}t^n\,\mathrm{d}\mu(t)\qquad(n\ge0).
\]
The supplied material raises moment-sequence questions about the fixed-occurrence coefficient sequences and about the weighted sequence \((g_{\tau,n}(q))_{n\ge0}\). However, the excerpt ends immediately after introducing the weighted coefficients: it supplies neither a complete conjecture about them nor a domain for \(q\). No definite universally quantified moment-sequence assertion can therefore be recovered.

The excerpt attributes to Bostan, Elvey Price, Guttmann and Maillard a Stieltjes-moment result for the coefficients of all pattern-avoiding-permutation generating functions described there as known. It also reports increasing Hankel determinants for \(\operatorname{Av}(1324)\) as the determinant size increases, motivating a conjecture that its counting sequence is a Stieltjes moment sequence. The determinant indexing is unspecified. For positive fixed occurrence counts, it reports failure of the Stieltjes moment property for every pattern and every \(r>0\) investigated; this is a qualification about the investigated cases, not a universal theorem. The weighted construction is attributed to Blitvic and Steingr\'{\i}msson.
\subsection{Short English statement}
For each of two unidentified length-four pattern classes, the proposed asymptotic says that prescribing a fixed number \(r\) of occurrences multiplies the avoidance count by approximately a constant times \(n^r\). The record also discusses Stieltjes moment properties of fixed-occurrence counts and occurrence-weighted counts. Missing class identities and an incomplete moment-sequence proposal prevent a fully specified statement. The supplied open-status assessment is dated curation and requires independent review.
\subsection{Sources}
\begin{itemize}
\item Supplied local UnsolvedMath registry record, permanent ID \(30005320\), title: \emph{Wilf Class Asymptotics with Fixed Pattern Occurrences}; exact source problem code: \texttt{catalogue entry}. The registry labels the problem open and the statement \texttt{corrected\_verified}; its verification note describes repairs to delimiters, notation, grammar and typography, rather than independent mathematical verification.
\item Supplied report citation: \emph{Enumerative Combinatorics} (2023), Oberwolfach Reports, DOI 10.4171/owr/2022/57, pp.\ 3215--3305. The registry gives proposed year 2022 and no named proposer. Report link, also supplied as a literature source labelled ``original OWR report'': \url{https://doi.org/10.4171/owr/2022/57}.
\item Supplied literature source labelled ``Conway--Guttmann primary article'': \url{https://www.combinatorics.org/ojs/index.php/eljc/article/download/v32i1p3/pdf}. The supplied assessment describes a 2025 study in which the asymptotic for both intended length-four Wilf classes is explicitly conjectural. No article title or fuller bibliographic details are supplied.
\item Literature check date: \texttt{2026-08-22}. The supplied assessment states that the asymptotic remains open. This is supplied dated curation derived from the OWR material, not an independently verified literature assessment; time-sensitive status and cited claims require independent verification.
\item The excerpt attributes the avoidance moment-sequence discussion to Bostan, Elvey Price, Guttmann and Maillard, and the weighted construction to Blitvic and Steingr\'{\i}msson. Separate authoritative citations for these attributed contributions are not supplied.
\end{itemize}
\end{document}
